% =========================================================================
% Formato y diseño del documento. PENDIENTE: se aplicará en la fase de
% diseño. Los requisitos ya conocidos (guía de TFG de la UPV/EHU) quedan
% anotados aquí para no perderlos:
%
%   - Márgenes: 2,5 cm por los cuatro lados        -> paquete geometry
%   - Interlineado sencillo; 6 pt entre párrafos   -> parskip / setspace
%   - Tipo de letra: Calibri                       -> con pdflatex no existe
%     Calibri; opciones: paquete `carlito` (clon métricamente compatible)
%     o compilar con XeLaTeX/LuaLaTeX + fontspec usando la Calibri de
%     Windows. Decidir en la fase de diseño.
%   - Títulos: capítulo MAYÚSCULA 14 pt negrita; apartado 1er nivel
%     MAYÚSCULA 12 pt negrita; 2º nivel minúscula 12 pt negrita; 3er nivel
%     minúscula 12 pt cursiva                      -> paquete titlesec
%   - Numeración de apartados: 1., 1.1., 1.1.1.    [ya es el defecto]
%   - Tablas/figuras: doble numeración capítulo.orden  [report ya lo hace]
%   - Impresión a doble cara                       -> \documentclass[..., twoside]
%   - Extensión objetivo: 40 páginas (máx. 60 excepcional)
%   - Identidad visual UPV/EHU en la portada (logos)
%
% Ejemplos de lo que vivirá aquí (dejar comentado de momento):
% \usepackage[a4paper, margin=2.5cm]{geometry}
% \usepackage{parskip}
% \usepackage{titlesec}
